\chapter*{ABSTRACT}
\addcontentsline{toc}{chapter}{ABSTRACT}

\begin{center}\textbf{\MakeUppercase{\ProjectTitle}}\end{center}

% 250-500 words. Cover: (1) problem, (2) methods, (3) results, (4) implications.
% Write this LAST, after the rest of the report has stabilised. See GUIDELINES.md
% \S Abstract for the structure to follow.

Vector search, the task of finding the stored items whose embeddings lie closest
to a query, underlies modern retrieval systems, including cross-modal search that
retrieves images from a text description. Two quantum primitives, Grover's
algorithm for unstructured search and the swap test for state similarity, appear
well suited to this task. This project investigates, through direct measurement,
whether they are of practical use for cross-modal vector search at the scale and
on the hardware available today.

We built a single framework in which classical, quantum, and hybrid search
engines answer the same queries over the same CLIP embeddings, so their rankings
can be compared on identical terms. Seven engines were implemented behind a
common interface: three classical baselines (an exact scan, an exact FAISS index,
and an approximate Hierarchical Navigable Small World graph), three quantum
engines (a swap-test estimator and two variants of Grover's algorithm), and a
hybrid that uses the classical index to shortlist candidates and the swap test to
rerank them. We benchmarked every engine over a controlled twenty-image
cross-modal dataset across a range of embedding dimensions and measurement shot
counts, recording retrieval accuracy as Mean Reciprocal Rank alongside
hardware-independent cost measures such as oracle calls, qubit count, and circuit
depth. The quantum engines ran on a classical simulator, with a single small
validation run on real IBM quantum hardware.

The results show that retrieval accuracy was governed by the embedding dimension
rather than by whether an engine was quantum. The swap test reproduced the
classical ranking accuracy once given a sufficient measurement budget, matching it
exactly at the highest dimension, and our Grover implementation reproduced its
theoretical oracle-call count exactly. No engine that performed its search
quantumly outperformed the classical graph-based baseline, however, and the
hybrid engine matched that baseline only because the classical index performed the
search. The underlying reason is structural: vector retrieval is dominated by
reading the dataset, a data-bound task, whereas quantum advantage concentrates on
compute-bound problems over small inputs. We conclude that cross-modal vector
search is not a natural fit for quantum methods, even setting aside the maturity
of current hardware, and that the apparent speed-up of quantum search rests on a
data-loading cost that the headline figures omit. The contribution of this work is
an honest, reproducible, like-for-like comparison and an explicit account of where
the quantum-classical gap actually lies.

\vspace{1em}
\noindent\textbf{Keywords:} \emph{quantum computing, vector search, similarity search, swap test, Grover's algorithm, CLIP, benchmarking}

\clearpage
